% ===== 原文 pp.19–20（§3.2 正则纤维续；3.2.1；分裂实形式） =====

% [原文 p.19]
表示 $b$ 处的纤维。纤维 $Y_b$ 的 $n$ 次同调在 $B$ 上构成一个局部平凡的向量丛，我们把它记作
$\mathcal{H}_n(B)$。该丛带有一个典范平坦联络，即 \emph{Gauss-Manin 联络}。

为定义这一联络，我们把 $\mathcal{H}_n(B)$ 在邻近点 $b_1,b_2\in B$ 处的纤维等同起来，即
$H_n(Y_{b_1})$ 与 $H_n(Y_{b_2})$。考虑包含 $b_1$ 与 $b_2$ 的可缩开集 $N\subset B$。纤维的包含
映射 $Y_{b_1}\hookrightarrow p^{-1}(N)$ 与 $Y_{b_2}\hookrightarrow p^{-1}(N)$ 都是同伦等价，
于是我们得到可缩开集 $N$ 中某点处纤维的同调与 $H_n(p^{-1}(N))$ 之间的同构：
\[
H_n(Y_{b_1})\cong H_n(p^{-1}(N))\cong H_n(Y_{b_2}).
\]

这意味着 $B$ 上的向量丛 $\mathcal{H}_n(B)$ 有一个平坦联络，即 Gauss-Manin 联络。
$\mathcal{H}_n(B)$ 的\emph{单值性}就是该联络的和乐，即一个同态
$\pi_1(B)\to\mathrm{Aut}(H_n(Y))$，也就是 $\pi_1(B)$ 在 $H_n(Y)$ 上的一个作用。把这些结果
应用于纤维化 $\mathcal{M}_{\mathrm{reg}}\to\mathcal{A}_{\mathrm{reg}}$，便知
$h:\mathcal{M}_{\mathrm{reg}}\to\mathcal{A}_{\mathrm{reg}}$ 的纤维上同调上的 Gauss-Manin
联络定义了 Hitchin 纤维化的单值作用。这一单值作用已对某些群作过研究，见 [12,16,88,89]。
% [Tutor Note: p.19 原文 "This monodromy action as been studied for some groups" 中 "as" 应为
%  "has"（原文笔误）；另注意此处 $\mathcal{M}_{\mathrm{reg}}\to\mathcal{A}_{\mathrm{reg}}$ 用的
%  是"正则纤维/正则轨迹"上的纤维化：只有去掉判别式轨迹之后才是局部平凡的，单值性才有定义。]

\para{3.2.1.\ Hitchin 纤维化与 $\Sl(n,\C)$-Higgs 丛}

为给出所需条件的一个例子，我们考虑 $G_{\C}=\Sl(n,\C)$。此时，由定义 2.9，一个
$\Sl(n,\C)$-Higgs 丛就是一个经典 Higgs 丛 $(E,\Phi)$，其中秩 $n$ 的向量丛 $E$ 有平凡行列式，
且 Higgs 场的迹为零。$\fsl(n,\C)$ 上不变多项式的一组基由迹为零的矩阵 $A\in\fsl(n,\C)$ 的特征
多项式的系数给出，于是由 Higgs 场定义的谱曲线 $\pi:S\to\Sigma$ 的方程为取 $a_1=0$ 的式 (3.5)。
此时 Hitchin 纤维化的一般纤维由 $\Jac(S)$ 中满足 $\pi_*L=E$ 且 $\Lambda^n\pi_*L$ 平凡的线丛
$L$ 构成的子集给出：$\Sl(n,\C)$ Hitchin 纤维化的一般纤维双全纯等价于 $\Prym(S,\Sigma)$，其中
$S$ 是按式 (3.5) 定义、系数 $a_1=0$ 的谱曲线。

为了体会由 Hitchin 纤维化产生的部分几何，我们接下来考虑 $n=2$ 的情形，并设 $\mathcal{A}$ 为
$\Sl(2,\C)$ 的 Hitchin 底，$\mathcal{M}$ 为 $\Sl(2,\C)$-Higgs 丛的模空间。对任意
$\omega\in\mathcal{A}-\{0\}$，[52, Theorem 8.1] 证明了纤维 $\mathcal{M}_\omega$ 是连通的。
对 $\mathcal{M}$ 中任一 $(E,\Phi)$ 的同构类，可以考虑其特征多项式的零点集
\[
\det(\Phi-\eta I)=\eta^2+\omega=0,
\]
其中 $\omega=\det(\Phi)\in\mathcal{A}$。这在 $K$ 的全空间 $X$ 中定义了一条谱曲线
$\rho:S\to\Sigma$，其中 $\eta$ 是 $K$ 的拉回丛的重言截面。我们用 $\mathcal{M}_{\mathrm{reg}}$
记 Hitchin 映射 $h$ 的正则纤维，并令 $\mathcal{A}_{\mathrm{reg}}$ 为底的正则轨迹，它由具有单

% [原文 p.20]
零点的二次微分给出。注意：在正则轨迹 $\mathcal{A}_{\mathrm{reg}}$ 上曲线 $S$ 是非奇异的，且
分歧点由 $S$ 与零截面的交给出。曲线 $S$ 具有一个自然对合 $\sigma(\eta)=-\sigma$，因此可以
把 Prym 簇 $\Prym(S,\Sigma)$ 定义为满足
\[
\sigma^*M\cong M^*
\]
的线丛 $M\in\Jac(S)$ 构成的集合。
% [Tutor Note: p.20 原文 "The curve S has a natural involution σ(η) = −σ" 中右端必为笔误，应为
%  σ(η) = −η：对合作用在 $K$ 的纤维坐标 $\eta$ 上，$S$ 的方程 $\eta^2+\omega=0$ 在
%  η ↦ −η 下不变（而 −σ 与 η 不是同一类对象，无法比较）；紧随其后的 Prym 定义
%  σ*M ≅ M* 也正依赖于这个 "η ↦ −η" 的二重覆盖对合。正文照原文排 "−σ"。]

特别地，这个定义与 $n>2$ 时借助范数映射（Norm map）给出的定义是相容的。

由于 $S$ 是 $K$ 的全空间 $X$ 中的二重覆盖，$\Prym(S,\Sigma)$ 中平凡丛 $\mathcal{O}$ 的正像由
$\pi_*\mathcal{O}=\mathcal{O}\oplus K^{-1}$ 给出（例如见 [18, Remark 3.1]）。此外，由自然
对合 $\sigma$，$\pi_*\mathcal{O}$ 的截面可以分成不变截面与反不变截面。由于 $\mathcal{O}$
对应于不变截面，且 $\Lambda^2\pi_*\mathcal{O}\cong K^{-1}$，故必有
$\pi_*\mathcal{O}=\mathcal{O}\oplus K^{-1}$。

如式 (3.3) 中引入的 Hitchin 截面，可以自然地由 $S$ 上的线丛 $L=\pi^*K^{1/2}$ 得到，对此有
\[
\pi_*\pi^*K^{1/2}=K^{1/2}\otimes\pi_*\mathcal{O}=K^{1/2}\oplus K^{-1/2}.
\]

因此，线丛 $\mathcal{O}\in\Prym(S,\Sigma)$ 有一个相关联的 Higgs 丛，它由
$(K^{1/2}\oplus K^{-1/2},\Phi_\omega)$ 给出，其中 Higgs 场 $\Phi_\omega$ 按例 2.5 与
「Example newexample」定义为
\[
\Phi_\omega=\begin{pmatrix} 0 & \omega\\ 1 & 0\end{pmatrix},
\qquad \omega\in H^2(\Sigma,K^2).
\]
% [Tutor Note: p.20 原文 "defined as in Example 2.5 and Example newexample by" 中
%  "Example newexample" 是原稿中未解析的交叉引用（源码中的标签即为 newexample），并非某个
%  真实存在的例子；本译本照原样保留该串。从上下文看，它应指例 3.6 中给出的那个局部形式
%  Φ = (0 w; 1 0)，与例 2.5 的 Φ_ω 同型。]
% [Tutor Note: p.20 同一式右端写 "for ω ∈ H^2(Σ, K^2)"，其中上标 2 应为 0：ω 是二次微分，即
%  $K^2$ 的全纯截面，属于 $H^0(\Sigma,K^2)$；例 2.5（c01）与式 (3.5) 中的系数 $a_1$ 都取自
%  $H^0$。正文照原文排 $H^2$。]

对一般的 $G_{\C}$，纤维的描述可以借助 Cameral 覆盖 [39]（另见 [37]）得到，这与后面几节对
经典李群给出的描述是等价的。关于实 Higgs 丛的 Cameral 覆盖的全面描述，读者可参阅 [82,83]。

\para{分裂实形式。}

为研究分裂实形式 $\Sl(2,\R)$ 的 Higgs 丛，应当记住：它们的谱数据对应于给出 Hitchin 纤维化
一般纤维的阿贝尔簇 $\Prym(S,\Sigma)$ 中的二阶元。此外，此时单值性由
$\pi_1(\mathcal{A}_{\mathrm{reg}})$ 在 $H^1(\Prym(S,\Sigma),\Z)$ 上的作用生成；而这个空间
又可证明给出纤维的 2-挠点集 $P[2]$（这可推广到更高秩，读者可参阅 [16] 中更详细的研究）。

下面我们将看到，空间 $H^1(\Prym(S,\Sigma),\Z)$ 如何同构于一般纤维 $\Prym(S,\Sigma)$ 中的
2-挠点 $P[2]$。由于 $h$ 的一般纤维由阿贝尔簇 $\Prym(S,\Sigma)$ 给出，即由复向量空间 $V$
对某个格 $\wedge$ 的商给出，故有相应的同调群正合列
\[
\ldots\to\pi_1(\wedge)\to\pi_1(V)\to\pi_1(\Prym(S,\Sigma))\to\pi_0(\wedge)\to\ldots.
\]
% [Tutor Note: p.20 末把 $H^1(\Prym(S,\Sigma),\Z)$ 直接说成同构于 2-挠点集 $P[2]$，严格说应
%  经过 mod 2：$H_1(\Prym,\Z)\otimes\Z/2\cong\Prym[2]=P[2]$。秩即可验证——覆盖次数为 2 时
%  $\dim\Prym(S,\Sigma)=g_S-g$，故 $H^1(\Prym,\Z)\cong\Z^{2(g_S-g)}$ 有 $2(g_S-g)$ 个自由
%  生成元，而 $P[2]$ 有 $2^{2(g_S-g)}$ 个元素，两者基数不同，不可能同构。]
% [Tutor Note: 上述同调正合列可用来核对最后一句：$V$ 可缩（$\pi_1(V)=\pi_0(V)=0$）、格 $\wedge$
%  离散（$\pi_1(\wedge)=0$，但 $\pi_0(\wedge)=\wedge$），故该列给出
%  $\pi_1(\Prym(S,\Sigma))\cong\pi_0(\wedge)=\wedge$——即阿贝尔簇的基本群就是定义它的格；
%  再由 $H^1(\Prym,\Z)\cong\Hom(\wedge,\Z)$ 即得 $\dim H^1=\rk\wedge=2(g_S-g)$，与上一条批注
%  的秩一致（对秩 2 谱曲线 $g_S=4g-3$，得 $2(g_S-g)=6g-6=\dim_{\R}\mathcal{M}$，可交叉验算）。]
